\section{Preserving a prescribed denominator}\label{sec:prescribed}

Membership in $D_k$ requires the exact denominator $m$.
Lemma~\ref{lem:marked-cleanup}, used for counting, preserves only a
multiple of a specified odd integer, so it does not ensure that the
term $1/m$ survives. We instead reserve $1/m$ and use a greedy
construction that skips denominator $m$. It produces a distinct prefix
and makes the remaining sum smaller than every reserved reciprocal.
Any positive expansion of this remainder then avoids both the marker
and the prefix denominators.

We use this prefix twice: the uniform theorem first gives a short tail,
and a direct construction later gives the numerical bound in
Corollary~\ref{cor:prescribed}. Keeping the remainder unreduced retains
the divisors needed for that direct construction. In either case, a
separate padding argument will retain $m$ at every larger length.

\begin{lemma}[A greedy prefix avoiding one denominator]
\label{lem:marked-greedy-prefix}
Let $m\ge4$ and $T\ge2m^2$ be integers.
There are integers $2\le n_1<\cdots<n_j$, none equal to $m$,
and integers $q>0$ and $0\le R<2m$ such that
\[
  1=\frac1m+\sum_{i=1}^{j}\frac1{n_i}+\frac Rq,
  \qquad q=m\prod_{i=1}^{j}n_i,
  \qquad
  j<3+\log_2\log_2 T.
\]
Writing $q_0=m$ and $q_i=m\prod_{h=1}^{i}n_h$ for $1\le i\le j$,
we have $n_i\le q_{i-1}+1$ for every $1\le i\le j$.
If $R>0$, then
\[
  q\ge T,\qquad
  \frac Rq<\frac{2m}{T}\le\frac1m,
  \qquad \frac Rq<\frac1{n_i}\quad(1\le i\le j).
\]
In either case $q<T^2$.
\end{lemma}

\begin{proof}
Start with $(R_0,q_0)=(m-1,m)$ and $x_0=R_0/q_0$.
Whenever $R_i>0$ and $q_i<T$, set
\[
 a=\left\lceil\frac{q_i}{R_i}\right\rceil,
 \qquad
 n_{i+1}=\begin{cases}a,&a\ne m,\\m+1,&a=m,\end{cases}
 \qquad
 (R_{i+1},q_{i+1})=(n_{i+1}R_i-q_i,n_{i+1}q_i).
\]
Do not cancel common factors in this pair.  Thus
$x_{i+1}=R_{i+1}/q_{i+1}=x_i-1/n_{i+1}$.

For an ordinary step, writing $n=n_{i+1}=a$ gives
\[
  \frac1n\le x_i<\frac1{n-1},\qquad
  0\le R_{i+1}<R_i,\qquad
  0\le x_{i+1}<\frac1{n(n-1)}\le\frac1n.
\]
Moreover $n<1/x_i+1$, so
\[
 x_{i+1}<x_i-\frac{x_i}{1+x_i}
          =\frac{x_i^2}{1+x_i}<x_i^2.
\]
In the exceptional step $a=m$, we have
$(m-1)R_i<q_i\le mR_i$.  Consequently
\[
  0<R_i\le R_{i+1}<2R_i,
  \qquad
  0<x_{i+1}
   <\frac2{(m-1)(m+1)}<\frac1{m+1},
\]
where the last inequality uses $m\ge4$.
Every step therefore leaves a remainder smaller than the term just
subtracted.  The next denominator, if there is one, is strictly larger.
After the exceptional step all later denominators exceed $m+1$, so that
step can occur at most once.  Before it, $R_i\le m-1$; after it, the
numerator is less than $2(m-1)$ and decreases at every subsequent step.
Thus $0\le R_i<2m$ throughout.
Since $R_i\ge1$ in an active step, $a\le q_i$ and hence
$n_{i+1}\le q_i+1$.

All selected denominators are at least $2$, so $q_i$ at least doubles
at each step.  The procedure thus stops after finitely many steps.
Its first denominator is $2$, and
$x_1=(m-2)/(2m)<1/2$.
All later steps square the upper bound for the remainder, apart from
at most one step that still decreases it.  Therefore, whenever $i\ge2$
and the remainder is positive,
\[
 x_i\le 2^{-2^{i-2}},\qquad q_i\ge\frac1{x_i}
                 \ge2^{2^{i-2}}.
\]
If the final step has index $j\ge3$, its preceding state is active,
and hence
\[
 2^{2^{j-3}}\le q_{j-1}<T.
\]
This gives the asserted bound for $j$; it is immediate when $j\le2$.
The same preceding state satisfies
$q_j\le q_{j-1}(q_{j-1}+1)<T^2$, since $T$ is an integer.

The displayed decomposition and product formula follow by telescoping.
If $R_j=0$, the decomposition is complete.  Otherwise the stopping rule
gives $q_j\ge T$ and $R_j<2m$, whence
$R_j/q_j<2m/T\le1/m$.
At each earlier step the new remainder was smaller than $1/n_i$,
and subsequent steps only decrease it.  This proves all the claimed
strict inequalities.
\end{proof}

To see directly how the uniform theorem supplies a prescribed denominator,
take $T=2m^2$.  The lemma gives $j=O(\log\log m)$ and either a
complete marked expansion or a positive remainder $R/q$ with
\[
  1\le R<q,\qquad 2m^2\le q<4m^4.
\]
For sufficiently large $m$, the uniform upper bound in
Theorem~\ref{thm:main} applies to $R/q$, without requiring it to be
reduced.  It supplies a distinct expansion with
$O(\log\log q)=O(\log\log m)$ terms.  Each reciprocal in this
tail is at most its total $R/q$, which is strictly below $1/m$ and
every $1/n_i$.  The tail therefore avoids the marker and all prefix
denominators.  Together these terms give a distinct expansion of $1$
containing $1/m$, of length $O(\log\log m)$.

The same construction also supplies the finite exceptional values of
$m$, independently of the cutoff $T$.  For any fixed $m\ge4$, omit the condition
$q_i<T$ and continue until the remainder is zero.  The nonnegative
integer $R_i$ decreases strictly at every ordinary step, and there is
at most one exceptional step.  The process therefore terminates and
gives a distinct expansion of $1$ containing $1/m$.  For $m=2,3$, use
$1=1/2+1/3+1/6$.

The padding injection in Section~\ref{sec:counting} may remove the
denominator we wish to preserve. The following lemma supplies the
different property needed here: it increases the length and retains
any one prescribed denominator.

\begin{lemma}[Preserving a prescribed denominator while padding]
\label{lem:exact-marker-padding}
Let $r\ge3$ be an integer, and suppose that
\[
 1=\sum_{i=1}^{r}\frac1{n_i},\qquad
 1\le n_1<\cdots<n_r,
\]
where the denominators are integers. For every
$m\in\{n_1,\ldots,n_r\}$ there is a representation of $1$ by
exactly $r+1$ distinct positive unit fractions that still contains
the exact denominator $m$. Consequently $D_r\subseteq D_{r+1}$.
\end{lemma}

\begin{proof}
This is the padding lemma of van Doorn and Tang
\cite[Lemma~2.1]{vanDoornTang2026}; we include a proof.
Every denominator is at least two, since a term with denominator one
would already exhaust the sum. The identity
\[
 \frac1b=\frac1{b+1}+\frac1{b(b+1)}
\]
replaces one term by two with distinct denominators larger than $b$.
If $m\ne n_r$, apply it to $n_r$; the new denominators exceed all
the unchanged ones, and $m$ remains.

Suppose that $m=v=n_r$, and let $s=n_{r-1}$.
The same split of $s$ works unless $v=s+1$ or $v=s(s+1)$,
since every other unchanged denominator is smaller than $s$.
In either exceptional case, if $s=ab$ with integers $a,b\ge2$, use instead
\[
 \frac1s=\frac1{s+a}+\frac1{b(s+a)}.
\]
The inequalities
\[
 s+1<s+a<b(s+a)<s(s+1)
\]
show that both new denominators avoid either exceptional value of
$v$ and all the unchanged denominators below $s$.

Finally, neither exceptional case can occur when $s=p$ is prime.
There are at least three original terms, so $p>n_{r-2}\ge2$ and
in particular $p$ is odd. All denominators other than $p$ and $v$
are below $p$; write the sum of their reciprocals as $A/B$ with
$p\nmid B$. The two exceptional choices give respectively
\[
 1-\frac AB
 =\frac1p+\frac1{p+1}
 =\frac{2p+1}{p(p+1)},
 \qquad\text{or}\qquad
 1-\frac AB
 =\frac1p+\frac1{p(p+1)}
 =\frac{p+2}{p(p+1)}.
\]
After clearing denominators and reducing modulo $p$, these equations
give respectively $0\equiv B\pmod p$ and $0\equiv2B\pmod p$,
both impossible. Thus in every valid case an unmarked term can be
replaced by two distinct terms, preserving $m$ and increasing the
length by exactly one.
\end{proof}

\section{A quantitative prescribed-denominator bound}
\label{sec:prescribed-quantitative}

The uniform theorem gives the short marked expansion above with an
unspecified absolute length constant. We now prove the following
numerical estimate by constructing the tail directly.

\begin{proposition}\label{prop:marked-length}
For every $\varepsilon>0$ there is an integer $m_\varepsilon$ such
that every integer $m\ge m_\varepsilon$ occurs as an exact denominator
in a representation of $1$ by distinct positive unit fractions with
at most
\[
 \left(\frac{257}{\log 2}+\varepsilon\right)\log\log m
\]
terms.
\end{proposition}

The construction has two ingredients. A common integer $K_m$ supplies
short representations of small integers by positive rationals whose
numerators divide $K_m$. The unreduced denominator of the greedy
remainder has divisors at every multiplicative scale. Combining these
properties gives a short expansion of the tail: divisors of the greedy
denominator control how many pieces are needed, and the rational
representations supplied by $K_m$ bound the number of unit fractions
used for each piece.

\subsection{A common supply of rational divisors}

We construct a moderately sized integer whose divisors, after division by
positive integers, represent every integer in a prescribed initial interval
with a bounded number of summands.  The denominators of these rational
summands need not agree.  This freedom lets us first represent a positive
multiple of a prime and then divide by its integer multiplier.

\begin{lemma}\label{lem:rational-divisor-supply}
For every sufficiently large integer $m$, put
\[
 L=\log\log m,\qquad Y=m^4,\qquad
 \ell(v)=\left\lceil\frac{4\log v}{\log 2}\right\rceil
 \quad(v\geq2).
\]
Let $P(v)$ be the product of the first $\ell(v)$ primes, and define
\[
 K_m=P(Y)^2\,\lfloor L\rfloor!.
\]
Then $K_m>m$ and
\begin{equation}\label{eq:supply-size}
 \log K_m\leq
 \left(\frac{32}{\log 2}+o(1)\right)\log m\log\log m
 \qquad(m\longrightarrow\infty).
\end{equation}
Moreover, every integer $1\leq s\leq Y$ has a representation
\begin{equation}\label{eq:rational-supply}
 s=\sum_{i=1}^{h}\frac{e_i}{t_i},\qquad
 h\leq16,\quad e_i\mid K_m,\quad e_i,t_i\in\mathbb Z_{>0}.
\end{equation}
Repeated summands are allowed.
\end{lemma}

\begin{proof}
The prime number theorem \cite[p.~305, (1.1)]{Selberg1949} gives
\[
 \log P(v)=(1+o(1))\ell(v)\log\ell(v)
 \ll\log v\log\log v.
\]
Since $\ell(Y)=16\log m/\log 2+O(1)$, and
$\log(\lfloor L\rfloor!)=O(L\log L)$, this proves
\eqref{eq:supply-size}.  Also $P(Y)\geq2^{\ell(Y)}\geq Y^4$,
so $K_m>m$ for large $m$.

We next prove an assertion whose threshold is independent of $m$.
Fix a sufficiently large odd integer $u$, and let $A$ be the set of positive
divisors of $P(u)$.  Write $N=|A|=2^{\ell(u)}\geq u^4$.
We estimate exceptional primes by counting congruent pairs, the
mechanism underlying Gallagher's larger sieve \cite{Gallagher1971}.
The coarse count below suffices.
For each prime $p\leq u$, let $A_p$ be the image of $A$ in
$\mathbb F_p$.  Call $p$ exceptional when $|A_p|\leq p^{3/4}$.
If $c_a$ is the number of members of $A$ in residue class $a$, then
Cauchy--Schwarz shows that the number of ordered pairs of distinct members
of $A$ congruent modulo an exceptional prime is at least
\[
 \sum_{a\in\mathbb F_p}c_a^2-N
 \geq\frac{N^2}{|A_p|}-N
 \geq\frac{N^2}{2u^{3/4}}.
\]
For any two distinct members $a,b\in A$, the nonzero integer $|a-b|<P(u)$
has at most $\log P(u)/\log 2$ distinct prime divisors.
Counting the same pairs prime by prime therefore bounds the number $B(u)$
of exceptional primes by
\begin{equation}\label{eq:supply-exceptional}
 B(u)\leq \frac{2u^{3/4}\log P(u)}{\log 2}
 \ll u^{3/4}\log u\log\log u.
\end{equation}

We use the quantitative form of Vinogradov's three-prime theorem: there is
an absolute $c>0$ such that every sufficiently large odd integer $u$ has
at least $cu^2/(\log u)^3$ ordered representations as a sum of three
primes; see the statement in \cite[(1)--(2)]{Kumchev1997}
and the exposition in \cite[Theorem~1 and Section~3.1]{KumchevTolev2005}.
To make the uniformity in odd $u$ explicit, its singular series is
\[
 \mathfrak S(u)=
 \prod_{p\mid u}\left(1-\frac{1}{(p-1)^2}\right)
 \prod_{p\nmid u}\left(1+\frac{1}{(p-1)^3}\right).
\]
For odd $u$ its factor at $2$ is $2$, and
\[
 \mathfrak S(u)\geq
 2\prod_{p>2}\left(1-\frac{1}{(p-1)^2}\right)>0.
\]
The infinite product is positive because the sum of the subtracted
quantities converges.  Thus the asymptotic formula supplies an absolute
lower bound with an absolute threshold, with no dependence on the
factorization of $u$.

At most $3uB(u)$ ordered prime triples summing to $u$ contain an
exceptional prime: choose its position, its value, and one further entry;
the last entry is then fixed.  By \eqref{eq:supply-exceptional},
\[
 3uB(u)\ll u^{7/4}\log u\log\log u
       =o\left(\frac{u^2}{(\log u)^3}\right).
\]
Consequently every sufficiently large odd $u$ is the sum of three primes
that are not exceptional for this same set $A$.

It remains to represent each such prime using divisors of $P(u)^2$.
Fix one of them, say $p$, and put $H=A_p$ and $h=|H|>p^{3/4}$.
Let $U$ be the product of two independent uniformly chosen members of
$H$, and write $\mathrm e_p(z)=\exp(2\pi i z/p)$.
The following classical bilinear estimate is recorded in
\cite[Proposition~1.1]{GlibichukKonyagin2007}; we give its short proof
and the five-product consequence needed here.
For every nonzero $r\in\mathbb F_p$, orthogonality and
Cauchy--Schwarz give
\begin{align*}
 \left|\sum_{a,b\in H}\mathrm e_p(rab)\right|^2
 &\leq h\sum_{a\in\mathbb F_p}
       \left|\sum_{b\in H}\mathrm e_p(rab)\right|^2\\
 &=ph^2.
\end{align*}
After division by $h^4$, this says
\[
 \left|\mathbb E\,\mathrm e_p(rU)\right|
 \leq\frac{\sqrt p}{h}<p^{-1/4}.
\]
For five independent copies $U_1,\ldots,U_5$, Fourier inversion now yields
\begin{align*}
 \Pr(U_1+\cdots+U_5=0)
 &=\frac1p\left(1+\sum_{r\ne0}
                   (\mathbb E\,\mathrm e_p(rU))^5\right)\\
 &\geq\frac{1-(p-1)p^{-5/4}}p>0.
\end{align*}
Here the lower bound follows by bounding the absolute value of the sum
of the nonzero Fourier coefficients.

Choose pairs of residues witnessing this positive probability and lift
each residue to any positive member of $A$ representing it.  The five
integer products $e_1,\ldots,e_5$ are positive divisors of $P(u)^2$,
and their sum is divisible by $p$.  Hence
\[
 t=\frac{e_1+\cdots+e_5}{p}\in\mathbb Z_{>0},\qquad
 p=\sum_{i=1}^{5}\frac{e_i}{t}.
\]
The three chosen primes thus represent $u$ in fifteen summands of the
required kind.

For a sufficiently large integer $s$, use $u=s$ if $s$ is odd and
$u=s-1$ otherwise.  In the even case append the summand $1=1/1$.
This uses at most sixteen summands.  If $s\leq Y$, monotonicity of
$\ell$ implies $P(u)\mid P(Y)$, so every numerator constructed above
divides $K_m$.

Finally, all preceding thresholds concern $s$ alone and are absolute.
Choose an integer $S_0$ above them.  For all sufficiently large $m$ we
have $\lfloor L\rfloor\geq S_0$.  Each remaining integer
$1\leq s\leq S_0$ then divides $\lfloor L\rfloor!$ and hence $K_m$,
so it has the one-summand representation $s=s/1$.  This proves the
assertion simultaneously for every $1\leq s\leq Y$.
\end{proof}

\subsection{Grouping the remaining numerator}

The denominator retained by the greedy procedure has divisors at every
multiplicative scale.  We first record this property, then use it to reduce
the remaining numerator to a short list of integers within the range of
the rational-divisor supply.

\begin{lemma}\label{lem:marked-divisor-density}
Let $m\ge2$ be an integer.  Suppose that $q$ is obtained from $m$ by a
finite sequence of replacements $q\mapsto nq$, where at each replacement
$n$ is a positive integer satisfying $n\le q+1$.  Then, for every real
number $w$ with $1\le w\le q$, there is a positive divisor $d$ of $q$ such
that
\[
  \frac wm\le d\le w.
\]
\end{lemma}

\begin{proof}
For the initial value $q=m$, the divisor $1$ works throughout $[1,m]$.
Suppose the property holds for $q$, and put $q'=nq$.  If $1\le w\le q$,
use a divisor of $q$, which is also a divisor of $q'$.  If $n\le w\le nq$,
apply the property at $w/n$ and multiply the resulting divisor by $n$.
These two cases cover $[1,nq]$ unless $q<n$ and $q<w<n$.  In that gap
the divisor $q$ itself works: $q<w$ and
\[
  w<n\le q+1\le mq.
\]
This proves the induction, including the real points between consecutive
integers.
\end{proof}

\begin{lemma}\label{lem:marked-grouping}
Let $m\ge2$, $q\ge1$, $K_m\ge1$, and $B\ge1$ be integers.  Suppose that, for every
real $w\in[1,q]$, the integer $q$ has a divisor in $[w/m,w]$.  Suppose
also that each integer $s\in[1,m^4]$ can be written as a sum of at most
$B$ positive rational numbers $e/t$, where $e\mid K_m$ and $t$ is a
positive integer.  Then, for every integer $X$ with $1\le X\le q$, the
rational number $X/(qK_m)$ is a sum of at most
\[
  B\left(\frac{\log X}{2\log m}+2\right)
\]
positive unit fractions, with repetitions allowed.
\end{lemma}

\begin{proof}
Write $X_0=X$, and repeatedly remove groups from the remaining integer.
If the current integer $X_i$ exceeds $m^4$, then
$w=X_i/m^2$ lies in $[1,q]$.  Choose $d_i\mid q$ with
\[
  \frac{X_i}{m^3}\le d_i\le\frac{X_i}{m^2},
  \qquad s_i=\left\lfloor\frac{X_i}{d_i}\right\rfloor.
\]
Thus $m^2\le s_i\le m^3\le m^4$, and the integer remaining after the
group $d_i s_i$ is removed satisfies
\[
  0\le X_{i+1}=X_i-d_i s_i<d_i\le\frac{X_i}{m^2}.
\]
If instead $0<X_i\le m^4$, remove the final group with $d_i=1$ and
$s_i=X_i$.  Stop immediately if the remaining integer is zero.

Every grouping step with $X_i>m^4$ decreases the remainder by a factor
strictly greater than $m^2$.  There are therefore at most
$\lceil\log X/(2\log m)\rceil$ such steps, followed by at most one
final group.  The number $G$ of groups satisfies
\[
  G\le\left\lceil\frac{\log X}{2\log m}\right\rceil+1
    \le\frac{\log X}{2\log m}+2.
\]
This bound also covers $X=1$ and a zero remainder before the final group.

For each group, use the assumed representation
$s_i=\sum_{a=1}^{b_i}e_{i,a}/t_{i,a}$, with $b_i\le B$.
Dividing $X=\sum_i d_i s_i$ by $qK_m$ gives
\[
  \frac{X}{qK_m}
    =\sum_i\sum_{a=1}^{b_i}
       \frac{1}{(q/d_i)(K_m/e_{i,a})t_{i,a}}.
\]
Each displayed denominator is a positive integer.  This uses the two
divisibilities $d_i\mid q$ and $e_{i,a}\mid K_m$; no divisibility between
$t_{i,a}$ and $e_{i,a}$, or coprimality assumption, is needed.  The total
number of terms is at most $BG$.
\end{proof}

\begin{proof}[Proof of Proposition~\ref{prop:marked-length}]
Fix a sufficiently large integer $m$, put $L=\log\log m$, and take
$K_m$ from Lemma~\ref{lem:rational-divisor-supply}. Apply
Lemma~\ref{lem:marked-greedy-prefix} with $T=2mK_m$, which is an
integer at least $2m^2$. It gives
\[
 1=\frac1m+\sum_{i=1}^{j}\frac1{n_i}+\frac Rq,
 \qquad 0\le R<2m.
\]
The prefix is distinct and avoids $m$. By \eqref{eq:supply-size},
\[
 \log\log(2mK_m)\le L+O(\log L),
 \qquad
 j\le\frac{L}{\log2}+O(\log L)
   =\left(\frac1{\log2}+o(1)\right)L.
\]
If $R=0$, the expansion is complete and already satisfies the asserted
length bound. Suppose henceforth that $R>0$.

The denominator $q$ is retained without cancellation, beginning at $m$
and using multipliers at most the preceding denominator plus one.
Lemma~\ref{lem:marked-divisor-density} therefore applies to this $q$.
Moreover $q\ge2mK_m$ and $R<2m$, so the integer $X=RK_m$ satisfies
$1\le X<q$. Lemma~\ref{lem:marked-grouping}, with $B=16$ and the
supply from Lemma~\ref{lem:rational-divisor-supply}, represents
$X/(qK_m)=R/q$ as a sum of at most $16G$ unit fractions, where
\[
 G\le\frac{\log(RK_m)}{2\log m}+2
 \le\frac{\log K_m+\log(2m)}{2\log m}+2
 \le\left(\frac{16}{\log2}+o(1)\right)L.
\]
This count depends on $RK_m$, not on the potentially much larger $q$.

Apply Lemma~\ref{lem:distinct} only to this tail. It makes its
denominators distinct without changing its number of terms. The tail
sum is strictly below $1/m$ and every $1/n_i$ by
Lemma~\ref{lem:marked-greedy-prefix}. Every reciprocal in any positive
expansion of that sum is at most the sum itself. Thus all denominators
of the distinct tail exceed $m$ and every $n_i$, so adjoining it to the
reserved term and prefix produces a distinct expansion of $1$.
There is one reserved term, $j$ prefix terms, and at most sixteen terms
for each of the $G$ groups. Its total length is therefore at most
\[
 1+j+16G
 \le\left(\frac{1+16\cdot16}{\log2}+o(1)\right)L
 =\left(\frac{257}{\log2}+o(1)\right)\log\log m.
\]
All errors here tend to zero as $m$ tends to infinity through all
integers: the bound on $R$ removed any dependence on the chosen
stopping numerator, and the supply covers both parities. Absorbing
the error into any prescribed $\varepsilon>0$ proves the proposition.
\end{proof}

\subsection{Every sufficiently large prescribed length}

\begin{proof}[Proof of Corollary~\ref{cor:prescribed}]
Set $A=257/\log2$. Fix a real number $c$ with $0<c<1/A$, and
choose $\varepsilon>0$ such that $c(A+\varepsilon)<1$.
Proposition~\ref{prop:marked-length} supplies an integer $M\ge4$ such
that every integer $m\ge M$ has a distinct marked expansion with
at most $(A+\varepsilon)\log\log m$ terms. The construction in
Section~\ref{sec:prescribed} also gives a finite distinct marked
expansion for each $2\le m<M$. Fix one for each such $m$, and let
$k_*$ be the maximum of their finitely many lengths.

For every integer $k\ge k_*$ and every integer
$2\le m\le\exp(\exp(ck))$, the chosen expansion has at most $k$
terms. This is immediate for $m<M$; for $m\ge M$ its length is at
most
\[
 (A+\varepsilon)\log\log m
 \le c(A+\varepsilon)k<k.
\]
A distinct expansion of $1$ containing a denominator $m\ge2$ has
at least three terms: denominator one would exhaust the sum, and two
distinct denominators at least two contribute at most $1/2+1/3<1$.
Thus Lemma~\ref{lem:exact-marker-padding} may be iterated until the
length is exactly $k$, retaining the exact integer $m$. We have proved
\[
 \{2,\ldots,\lfloor\exp(\exp(ck))\rfloor\}\subseteq D_k.
\]
The first missing integer exceeds the real right endpoint, and hence
$v(k)\ge\exp(\exp(ck))$ for every sufficiently large integer $k$.
Because this holds for each fixed $c<\log2/257$, it gives
\[
 \liminf_{k\to\infty}\frac{\log\log v(k)}{k}
 \ge\frac{\log2}{257}.
\]
In particular $257/\log2<600$, so taking $c=1/600$ gives the stated
eventual lower bound $v(k)\ge\exp(\exp(k/600))$.

For the upper bound, the product recurrence in the proof of
Corollary~\ref{cor:counting} gives $n_i\le k^{2^{i-1}}$ in every
distinct exact-$k$ expansion of $1$. Every denominator in such an
expansion is therefore at most $k^{2^{k-1}}$. Consequently $D_k$ is
finite for every $k\ge1$, its complement in $\{2,3,\ldots\}$ is
nonempty, and
\[
 v(k)\le1+k^{2^{k-1}}.
\]
For $k\ge2$ this yields
$\log\log v(k)\le k\log2+\log\log k$, and therefore
\[
 \limsup_{k\to\infty}\frac{\log\log v(k)}{k}\le\log2.
\]
Together the two bounds imply $\log\log v(k)=\Theta(k)$.
\end{proof}

The lower slope above is the consequence of the length coefficient
in Proposition~\ref{prop:marked-length}. Its vanishing error does not
imply an eventual inequality with the endpoint $c=\log2/257$.
Neither that endpoint inequality nor a limiting or sharp slope is
asserted.

